% from ega4-2-fr.tex
\oldpage[IV-2]{24}
Reste à prouver (2.5.5.1)~; on a
$(\mathscr E\otimes_Y\mathscr F)_x
=\mathscr E_y\otimes_{\mathcal O_y}\mathscr F_x$
(\textbf{I}, 9.1.12)
et
$(\mathscr F_y)_x=\mathscr F_x/\mathfrak m_y\mathscr F_x$~; on est donc
ramené à prouver le

\begin{lemma}[2.5.5.2]
\label{IV.2.5.5.2-fr}
Soient $A$, $B$ deux anneaux locaux, $\rho:A\to B$ un homomorphisme local,
$\mathfrak m$ l'idéal maximal de $A$. Soient $M$ un $A$-module, $N$ un
$B$-module qui est un $A$-module fidèlement plat et tel que
$N/\mathfrak mN$ soit un $B$-module de longueur finie~; alors on a
\end{lemma}

\begin{env}[2.5.5.3]
\label{IV.2.5.5.3-fr}
\[
\operatorname{long}_B(M\otimes_A N)
=\operatorname{long}_A(M).
 \operatorname{long}_B(N/\mathfrak mN).
\]
\end{env}

Si $M$ est de longueur infinie, il en est de même de $M\otimes_A N$, car
pour toute suite strictement croissante de $n$ sous-modules $M_i$ de $M$
($1\leq i\leq n$), les $M_i\otimes_A N$ s'identifient à des sous-$B$-modules
de $M\otimes_A N$ deux à deux distincts~; comme
$N\ne\mathfrak mN$ puisque $N$ est un $A$-module fidèlement plat, la formule
(2.5.5.3) est vraie dans ce cas. Supposons donc $M$ de longueur finie. Si
$M=0$ les deux membres de (2.5.5.3) sont nuls, donc on peut supposer
$M\ne0$. Les $\mathfrak m^kM$ sont des sous-modules de $M$, donc de longueur
finie, et si $\mathfrak m^kM\ne0$, il résulte du lemme de Nakayama que
$\mathfrak m^{k+1}M\ne\mathfrak m^kM$~; donc il y a nécessairement un entier
$r$ tel que $\mathfrak m^rM=0$. Les
$\mathfrak m^kM\otimes_A N$ s'identifient alors à des sous-$B$-modules de
$M\otimes_A N$,
$(\mathfrak m^kM\otimes_A N)/(\mathfrak m^{k+1}M\otimes_A N)$
étant isomorphe à
$(\mathfrak m^kM/\mathfrak m^{k+1}M)\otimes_A N$,
donc aussi à
$(\mathfrak m^kM/\mathfrak m^{k+1}M)
\otimes_{A/\mathfrak m}(N/\mathfrak mN)$
$(0\leq k\leq r\setmin 1)$~;
la longueur de ce dernier, considéré comme $B$-module, est donc le produit de
$\operatorname{long}_B(N/\mathfrak mN)$ et du rang du
$(A/\mathfrak m)$-espace vectoriel
$\mathfrak m^kM/\mathfrak m^{k+1}M$, rang qui est aussi la longueur du
$A$-module $\mathfrak m^kM/\mathfrak m^{k+1}M$. Par addition (pour
$0\leq k\leq r\setmin 1$), on en déduit aussitôt la formule (2.5.5.3).

\begin{remark}[2.5.5.4]
\label{IV.2.5.5.4-fr}
On notera que lorsque $N$ est un $A$-module \emph{de type fini} dire que $N$
est un $A$-module fidèlement plat équivaut à dire que $N\ne0$ et que $N$ est
un $A$-module \emph{plat}~; en effet le lemme de Nakayama montre alors que
$\mathfrak mN\ne N$.
\end{remark}

\begin{lemma}[2.5.6]
\label{IV.2.5.6-fr}
Soient $B$ un anneau non nécessairement commutatif, $V$, $W$ deux
$B$-modules à gauche isomorphes, $C$ l'anneau
$\operatorname{End}_B(V)$, et soit
$M=\operatorname{Hom}_B(V,W)$, muni de sa structure canonique de
$C$-module à droite. Alors $M$ est un $C$-module isomorphe à $C_d$~; en
outre, pour tout $u\in M$, les conditions suivantes sont équivalentes~:
\begin{enumerate}
\item[\emph{a)}] $\{u\}$ est une base du $C$-module $M$.
\item[\emph{b)}] $u$ est un isomorphisme de $V$ sur $W$.
\end{enumerate}
\end{lemma}

Si $u$ est un isomorphisme de $V$ sur $W$, l'application
$v\leadsto u\circ v$ de $C$ dans $M$ est évidemment une bijection, donc
\emph{b)} entraîne \emph{a)}. Inversement, supposons que $\{u\}$ soit une
base du $C$-module $M$. Par hypothèse, il existe un isomorphisme $u'$ de $V$
sur $W$, et $\{u'\}$ est par suite une base de $M$~; il y a donc un élément
inversible $w$ de $C$ (c'est-à-dire un automorphisme de $V$) tel que
$u=u'\circ w$, ce qui entraîne que $u$ est un isomorphisme de $V$ sur $W$.

\begin{corollary}[2.5.7]
\label{IV.2.5.7-fr}
Les hypothèses sur $B$, $V$ et $W$ étant celles de \textup{(2.5.6)}, supposons en
outre vérifiée l'une des conditions suivantes~:
\begin{enumerate}
\item[\textup{(i)}] $V$ et $W$ sont des $B$-modules noethériens~;
\item[\textup{(ii)}] $V$ et $W$ sont des modules de présentation finie sur un
sous-anneau commutatif de $B$.
\end{enumerate}
Alors les conditions \emph{a)} et \emph{b)} de \textup{(2.5.6)} sont aussi
équivalentes aux suivantes~:
\begin{enumerate}
\item[\emph{a}$'$)] $u$ est un générateur du $C$-module $M$.
\item[\emph{b}$'$)] $u$ est un épimorphisme de $V$ sur $W$.
\end{enumerate}
\end{corollary}

On sait en effet qu'un épimorphisme d'un $A$-module $E$ sur lui-même est
\emph{bijectif} dans les deux cas suivants~: $1^\circ$ $E$ est un $A$-module
noethérien (Bourbaki, \emph{Alg.}, chap.~VIII, \S~2,
n\textsuperscript{o}~2, lemme~3)~; $2^\circ$ $A$ est commutatif et $E$ est
un $A$-module de présentation finie (8.9.3)\footnote{Le lecteur pourra
vérifier que (2.5.7) et (2.5.8) ne sont pas utilisés avant le \S~9.}~; donc
\emph{b)} et \emph{b}$'$) sont équivalentes. D'autre part, si $u$ engendre
$M$
